\section{Economic value of common-readout qualification}\label{sec:study84}
\providecommand{\FullMenuProblems}{976}
\providecommand{\JointProblems}{1572}
\providecommand{\ExactInfeasible}{104}
\providecommand{\IndependentInequalities}{16548}
\providecommand{\CoarseReLUBefore}{[-0.010033, 0.012680]}
\providecommand{\CoarseReLUAfter}{[-0.010023, 0.008293]}
\providecommand{\FineBasisAfter}{[0.000549, 0.002455]}
\providecommand{\BernsteinResolutionBefore}{[-0.000149, 0.005497]}
\providecommand{\BernsteinResolutionAfter}{[0.000210, 0.005089]}
\providecommand{\BernsteinRefinementFee}{0.000210}
\providecommand{\CompleteScientificSeconds}{104.699}

The new calculation refines the economic accuracy of the original attained NBO controllers. It does not fit another controller or change a conditional cost observation. The four calibrated readouts, their frozen labels, terminal actors, robust 33-action menus, acquisition grids and original continuous laws are those of the preceding factorial investment study. The original actionwise integrations and simultaneous attained-policy intervals remain the paid inputs.

\subsection{A complete, fixed qualification design}
At acquisition sizes eight and sixteen, every alternative admitted by the reference-only barrier at excess objective at most $2^{-12}$ receives a full-menu qualification problem. This gives \FullMenuProblems\ problems: 832 for coarse ReLU, nine for coarse Bernstein, 78 for fine ReLU and 57 for fine Bernstein. Each event includes the other 32 feasible menu actions, with a common coefficient vector and the same objective.

The joint-event design permits at most 512 pairs per model. When at most 128 alternatives remain at the largest allowance, all remaining alternatives seed the pair catalogue. Otherwise the minimum-lower and maximum-upper conditional-cost alternatives in each cell seed it. Cross-cell pairs are ranked by the product of their maximum absolute cost endpoints, with lexicographic tie handling, and the first 512 are retained. The complete roster contains \JointProblems\ pairs; the small Bernstein model contributes all 36 of its possible pairs. This rule is deterministic postprocessing of an already covered cost event, not a newly randomized validation design. A small implementation pilot preceded the source lock and is disclosed separately.

All six declared allowances are evaluated: zero and $2^{-k}$ for $k=16,15,14,13,12$. The output preserves singleton bounds, pair bounds, positive box separations, retained alternatives, active conflicts and both economic duals. Every one of the \ExactInfeasible\ infeasible optimization statuses has an independent positive rational separation certificate. A numerical status alone is never used to remove an event.

\begin{table}[htbp]
\centering
\caption{Original-law cost deviations at excess qualification objective $2^{-12}$}\label{tab:compatibility84}
{\small\setlength{\tabcolsep}{3pt}
\begin{tabular}{lrrrrcc}
\hline
Model & Single & Joint & Allowed & Cuts & Actionwise & Compatible \\
\hline
R8 & 832 & 512 & 691 & 351 & [-0.010033, 0.012680] & [-0.010023, 0.008293] \\
B8 & 9 & 36 & 9 & 14 & [-0.001080, 0.000900] & [-0.001032, 0.000792] \\
R16 & 78 & 512 & 61 & 140 & [-0.001168, 0.001291] & [-0.001009, 0.001053] \\
B16 & 57 & 512 & 56 & 204 & [-0.001311, 0.001311] & [-0.001012, 0.001000] \\
\hline
\end{tabular}}
\par\smallskip
{\footnotesize R and B denote calibrated ReLU and Bernstein readouts; the suffix is the acquisition size in each state coordinate. Bounds concern each optimizer class relative to its attained controller. Single and joint are qualification-problem counts; allowed and cuts apply at the displayed tolerance. Both interval endpoints are rounded outwards. Every tolerance, including unchanged bounds, appears in the appendix.}
\end{table}


The full-menu test gives a strict contraction of the inherited interval in eight of the 24 model--allowance records. Shared-coefficient restrictions give a further strict contraction in fifteen records. Some smaller-tolerance records already have no alternative action or no useful selected conflict and remain unchanged. The table and appendix retain both outcomes. For coarse ReLU at the largest allowance, the inherited deviation interval \CoarseReLUBefore\ becomes \CoarseReLUAfter. This is a reduction in the uncertainty about a class of possible policies, not another measured gain of the attained reference policy.

\subsection{A new classwide refinement decision}
All six pairwise comparisons at the largest allowance appear in Table~\ref{tab:comparisons84}. They add the compatible deviation bounds to the same attained-cost intervals used in the preceding study. The result is simultaneous over the declared coefficient classes, rather than a comparison of two numerical certificate values.

\begin{table}[htbp]
\centering
\caption{All paired optimizer-class comparisons at excess objective $2^{-12}$}\label{tab:comparisons84}
{\small\setlength{\tabcolsep}{3pt}
\begin{tabular}{lccr}
\hline
First minus second & Actionwise interval & Compatible interval & Sign \\
\hline
R8--B8 & [-0.007360, 0.018447] & [-0.007242, 0.014012] & unresolved \\
R8--R16 & [-0.006468, 0.019812] & [-0.006221, 0.015267] & unresolved \\
R8--B16 & [-0.005006, 0.021441] & [-0.004685, 0.016754] & unresolved \\
B8--R16 & [-0.001622, 0.003879] & [-0.001336, 0.003612] & unresolved \\
B8--B16 & [-0.000149, 0.005497] & [0.000210, 0.005089] & positive \\
R16--B16 & [-0.001523, 0.004614] & [-0.001054, 0.004077] & unresolved \\
\hline
\end{tabular}}
\par\smallskip
{\footnotesize Intervals use the same R80 simultaneous attained-controller cost event and the same original economic law. Positive means that the first controller class has higher cost throughout the certified pair of classes. An interval containing zero remains unresolved. These rows add no fitting attempt or independent observation.}
\end{table}


For coarse versus fine Bernstein, the previous optimizer-class interval at allowance $2^{-12}$ is \BernsteinResolutionBefore. It does not identify a signed ordering. The common-readout interval is \BernsteinResolutionAfter, whose lower endpoint is strictly positive. Thus the finer Bernstein controller has lower original-law cost for every pair of readouts in the two stated sublevels. Every additional normalized refinement charge below \BernsteinRefinementFee\ is supported by this classwide saving. This conclusion holds throughout smaller allowances by containment of the qualification classes, not by interpolation between the six printed tolerances.

The refinement decision is distinct from the matched basis comparison. At allowance $2^{-16}$ the fine ReLU-minus-Bernstein interval is \FineBasisAfter. It retains the identified Bernstein advantage under calibrated qualification while narrowing the preceding actionwise band. The earlier legacy qualification advantage of ReLU remains part of the factorial evidence. The paper therefore identifies the qualification rule and its verified optimization accuracy as economically relevant, rather than choosing a basis label and suppressing the opposite ranking.

The complete catalogue also retains coarse-ReLU comparisons that remain unresolved, and larger-allowance fine-basis comparisons whose intervals cross zero. An interval strictly inside the declared $1/256$ margin supports the corresponding equivalence statement; crossing zero alone does not. The new support calculation supplies no population fitting-success probability and does not convert a normalized cost unit into an estimated currency amount. The original priced-selection theorem can instead use the verified interval together with a separately supplied loss scale, deployment count and refinement charge.

\subsection{Complete incremental work and independent arithmetic}
The recorded scientific interpreter, including imports, reads, all qualification and assignment problems, durable outputs and shutdown, takes \CompleteScientificSeconds\ seconds on the stated ARM runner. Singleton, joint-event and assignment phases are recorded separately. These measurements price the additional qualification calculation. They are neither a cold-start cost for rebuilding the original controllers nor a matched ReLU-versus-Bernstein timing experiment. Original fitting and integration clocks remain unchanged and must be added when those inputs have not already been purchased.

A separate rational implementation reconstructs the feature antiderivatives, all winner comparisons, qualification supports, pair-event separations, conflict rosters and \IndependentInequalities\ assignment inequalities. It checks all 2,548 qualification certificates and all 36 policy-class comparisons without invoking the production optimizer. Original R77/R80 continuous-bin and actionwise kernel verification is separately replayed in the full publication build. The distinction between new exact certificate verification, inherited full-bin re-integration and independent cost sampling is explicit: this study adds no new fitting attempt, integration bin or independent cost row.

The economic improvement is consequently a more informative decision from the original NBO continuation and its previously qualified evidence. Full-menu winner regions remove alternatives that cannot be selected; common-coefficient restrictions remove jointly impossible selections; conditional integration prices what remains. Each step has its own proof, numerical certificate and work account. The original broader results and the difficult vector-horizon observations retain their separate assumptions and evidence.
